
\cleardoublepage
\newpage
\thispagestyle{empty}
\phantomsection
\addtotoc{Summary}

\begin{center}\huge \textbf{Summary}\end{center}

\vfill

\begin{flushleft}
\large

Security operations centres face a structural gap between detection and response. SIEM
platforms collect, correlate, and surface alerts---but what happens next remains a manual
decision. Commercial SOAR tools execute pre-scripted playbooks, which break down on
novel attack patterns and cannot reason about threat severity or asset criticality.

This thesis fills that gap with an intelligent reaction layer that sits downstream of a
SIEM and upstream of any action. The system takes a normalised alert, enriches it with
asset, identity, and threat-intelligence context, produces an explainable threat
assessment, generates a structured response plan, and passes every proposed action
through a deterministic policy engine before execution. All decisions are recorded in a
structured audit log linking each action to its evidence.

The architecture is a modular monolith: a C\# core engine enforces the orchestration
pipeline and safety constraints, while a Python intelligence service provides scoring
and planning via HTTP. The threat scorer combines TF-IDF logistic regression with
contextual enrichment boosts. The planner is a custom 270K-parameter multi-head neural
network that classifies strategy, predicts action sets, and regresses priority---
resolving entity parameters deterministically rather than generating them, eliminating
the hallucination risk that made LLM-based planning unsafe for live execution.

The planner achieves 96.4\% strategy accuracy and 97.2\% action F1 in end-to-end
pipeline evaluation---matching a 250-million-parameter fine-tuned seq2seq model at
under 5 milliseconds inference time. The policy engine guarantees that no model error
can cause a hard-denied action to execute. The full system runs on a single development
machine and exposes every stage of the decision pipeline for analyst review.

\end{flushleft}
